% Section 1 -- Introduction. Author's draft of 2026-09-05, cite keys mapped to refs.bib.
% Framing rule 1 (CLEAN START) is structural here: the motivation is the production deployments and the literature,
% never this project's own history. The closing sentence deliberately repeats the abstract's.
\section{Introduction}
\label{sec:intro}

Activation probes---small linear classifiers applied to a language model's hidden states---occupy an attractive point
in the safety-tooling design space. The activations they read are computed during inference regardless; the probes
themselves add microseconds of latency and kilobytes of memory; and unlike guard models, they scale to every request
at negligible marginal cost. This promise is no longer speculative: in early 2026, two frontier laboratories
independently reported production deployments of probe-first monitoring cascades, in which lightweight probes screen
all traffic and expensive classifiers examine only flagged cases \cite{cunningham2026classifiers, gdm2026probes}.

Production systems, however, train their probes on data at a scale most practitioners do not have. This paper asks how
far the approach stretches in the opposite regime---the minimal-data regime, where probes are trained on tens of
contrastive examples authored in an afternoon. The question is practically urgent: the minimal-data recipe is cheap
enough that enterprises, open-source projects, and researchers will attempt it, and recent work reports strong results
from exactly such recipes \cite{han2025safeswitch, conceptguard2025, mckenzie2025probes}. It is also scientifically
unsettled, because the evaluations behind those results differ from one another in choices that are rarely treated as
experimental variables at all: which token position is read, whether the prompt is formatted through the model's chat
template, which tensor a serving framework's hook actually returns, and what the benign side of the evaluation
contains.

We evaluate the minimal-data regime systematically: three safety tasks---harmful-content detection, prompt-injection
detection, and policy compliance---across seven open models from three architecture families, under a pre-registered
protocol with controls that, to our knowledge, no prior activation-probe study applies together: topic-matched benign
sets, length stratification, encoded-benign controls, refusal separation, lexical baselines, and leakage-grouped
splits.

Our central finding is that reported outcomes in this regime are dominated by two factors that most papers do not
report: extraction convention and evaluation design. The same probes, on the same models, move from apparent failure
to strong performance---or the reverse---when either factor changes. Concretely: probes that flag $30$--$94\,\%$ of
topic-matched benign prompts at their shipped thresholds become accurate detectors (length-stratified AUC $0.92$--$0.98$, exceeding tuned lexical
baselines by $+0.21$ to $+0.27$) when they are fitted with on-topic negatives and read through the model's chat
template rather than as bare text; and
probes that the conventional final-token read leaves at $0.58$--$0.85$ reach $0.94$--$0.97$ AUC when the activation is
instead averaged over the tool-response span, on the four models that clear a pre-registered pass threshold and on
injections never seen in training.
Neither repair changes the probe or its training data. Conversely, rigorous evaluation retires claims that weaker
evaluation supports: a proposed training method whose estimator we show to be algebraically identical to its baseline
and whose data construction we measure at no effect;
obfuscation robustness that can be manufactured at will by training on encodings; and a compliance-classification task
whose short-text evaluations, we show, cannot certify a semantic claim for any classifier---probe or
otherwise---because directive force in short text is carried by the very lexical markers such evaluations must
control.

The frontier-scale evidence points the same direction: DeepMind's production report identifies distribution shift in
input length as the principal generalization failure for deployed probes \cite{gdm2026probes}---at production scale,
the same class of phenomenon we localize, at accessible scale, to specific and repairable conventions.

We make five contributions.

\begin{enumerate}[leftmargin=*]
\item \textbf{Two extraction-convention repairs.} Reading prompts through the model's chat template removes the topic
      confound in harmful-content probing, raising length-stratified AUC by $+0.30$ to $+0.54$ on all seven models
      for probes fitted with on-topic negatives;
      and reading the mean over the tool-response span rather than the final token raises prompt-injection detection
      from $0.58$--$0.85$ to $0.94$--$0.97$ AUC on the four models that clear the pre-registered pass threshold, on
      injections never seen in training (Section~\ref{sec:extraction}).
\item \textbf{A construct-validity result.} In short authored text, directive force is substantially lexically
      constituted, so short-text authored evaluation cannot certify a semantic claim for any compliance classifier;
      we show this by building three instruments that each fail in a different way (Section~\ref{sec:learn}).
\item \textbf{An anatomy of obfuscation failure.} Across $22$ blocks, two template modes, two probe constructions,
      two activation spaces and three read positions, no configuration separates encoded-harmful from
      encoded-benign prompts once length is controlled---even for encodings the model demonstrably decodes---while training on encoded examples
      manufactures near-perfect apparent detection at a near-perfect encoded-benign false-positive rate
      (Section~\ref{sec:obfuscation}).
\item \textbf{Five evaluation hazards of independent value.} Exploitable label artifacts in a standard
      prompt-injection benchmark; prompt length as a pervasive confound; a tensor-identity hazard in serving-framework
      activation hooks; the necessity of separating probe detections from the model's own refusals; and the
      invalidity of leave-one-out cross-validation on minimal-pair data (Sections~\ref{sec:setup},
      \ref{sec:protocol}, \ref{sec:tensor}).
\item \textbf{A pre-registered protocol stress-tested in use.} Committed thresholds, two distinct amendment types for
      when a pre-registered interpretation fails, a stop rule that treats over-performance as a warning, and three
      cases in which our own pre-registered designs proved unable to answer their own questions
      (Section~\ref{sec:protocol}).
\end{enumerate}

Throughout, every number derives from committed, provenance-stamped results; the protocol, its pre-registered
thresholds, and the cases where our own pre-registered designs failed to answer their questions are reported as
findings in their own right. Minimal-data activation probes can work; whether they do is decided by conventions most
papers never report.
